\chapter{Kennismaken met de 45GS02-CPU}
\label{cha:introducing-the-45gs02-cpu}

Aan dit hoofdstuk wordt nog gewerkt.

\begin{itemize}
\item Hoe werkt een CPU\item Kenmerken van de 45GS02
\end{itemize}

\section{Werken met de monitor}

\begin{itemize}
    \item wat is een machinetaalmonitor\item MONITOR-commando\item registers bekijken\item geheugen bekijken\item code disassembleren\item code assembleren\item een subroutine uitvoeren\item onderbreken en hervatten
\end{itemize}

\section{CPU-instructies}

\begin{itemize}
    \item Instructies\item Adresseermethoden
\end{itemize}

\begin{itemize}
    \item Registers en statusvlaggen\item Geheugen lezen en schrijven\item Berekeningen uitvoeren\item Vertakkingen en subroutines\item MAP en EOM
\end{itemize}

\section{CPU-registers}

\begin{itemize}
    \item Accumulatorregister\item X, Y, Z\item Programmateller\item Stackpointer\item Base-Page (B)\item Statusvlaggen; het statusregister\item MAP
\end{itemize}

\section{De Base-Page}

\begin{itemize}
    \item Doel: snel tijdelijk geheugen; pointers\item Zero-Page\item Het B-register\item Base-Page-adresseermethoden\item Andere Base-Page-bewerkingen\item KERNAL-overwegingen
\end{itemize}

\section{Indirecte adressering}


\section{Het virtuele 32-bits register}


\section{De stapel}

\begin{itemize}
    \item Waarden op de stapel plaatsen en van de stapel halen\item Subroutine-aanroepen\item 8-bits versus 16-bits stackpointer\item Stackrelatieve adresseermethode
\end{itemize}

\section{Interrupts}

\begin{itemize}
    \item Typen interrupts\item Interruptvectoren\item Interface voor interrupthandlers\item Een voorbeeld zonder de KERNAL\item De KERNAL-IRQ uitbreiden\item Opmerking over het MAP-register
\end{itemize}

\section{Instructies coderen}

\begin{itemize}
    \item Opcodes en operanden\item 32-bits indirecte adresseermethode\item Instructies voor het virtuele 32-bits register
\end{itemize}

\section{De 6502-familie}

\begin{itemize}
    \item 6502 / 6510\item 65C02\item 65CE02 (``4502'') / 4510\item 45GS02
\end{itemize}
